\section{引言：把“宏观规律”化为一个判定问题}
\label{sec:introduction}

一个宏观描述之所以是一条可运行的规律，不只因为它使用较少的状态，还因为它在遗忘微观差异以后仍能预测自身的下一步。本文据此把“宏观规律何时存在”限定为一个明确的问题：同一表示能否同时保留指定任务、显著压缩操作上有意义的微观状态，并使压缩后的动力学近似闭合？

具体地，对一列规模增长的有限 Markov 系统，我们问：

\begin{quote}
在任务、微观探针、允许动力学和一步误差度量预先固定后，是否存在一列任务保真表示，使其相对状态数与最优闭合误差同时趋于零？
\end{quote}

这是一项条件判定，而不是“复杂系统必有简单规律”的存在宣言。本文采用最低限度的操作主义立场：可表达、可检验的规律依赖于对经验情形的区分、标记和概率关系；因此，数学模型必须先说明哪些区分属于当前任务和系统边界。该立场解释了任务与微观探针为何必须先于定理给定，但不替代后文需要验证的动力学条件。

\subsection{两个必须先排除的平凡化}

第一种平凡化来自\emph{无意义填充}。若给每个状态附加一个不影响任务和动力学的复制标签，原始状态数可以任意放大，任何固定表示都会显得高度压缩。为此，本文先由微观探针和允许动力学构造操作性微观空间 $B$，删除严格不可区分的复制；后续压缩只相对于 $|B|$ 计算，而且候选宏观表示必须是 $B$ 的任务保真信息商。

第二种平凡化来自\emph{只检查一半条件}。常值映射状态最少，却通常丢失任务；恒等映射永远精确闭合，却没有简化。即使某个候选被预先假定为低复杂度且近似闭合，从该假设推出“宏观规律存在”也只是在改写假设。本文因此先求单个有限系统的最粗精确任务商和完整的复杂度--误差前沿，再用不含预设宏观转移核的量 $\Gamma_N$ 判断这条前沿能否随系统规模逼近“相对复杂度为零、闭合误差为零”的原点。

\subsection{本文建立的结果}

全文的结论分为三个层次。

\begin{enumerate}
  \item \textbf{有限规模的规范基准。} 操作性约化不改变后续商的最坏行闭合误差。在约化空间上，对任意非空允许核族，都存在重标记意义下唯一的最粗任务保真精确闭合商；允许误差后，每个状态预算对应的最小闭合误差仍实际达到，并形成单调的复杂度--闭合前沿。
  \item \textbf{增长系统族的充要判据。} 纤维预测直径直接比较将被合并的状态对下一宏观块的预测。由半径--直径夹逼，$\Gamma_N\to0$ 当且仅当存在同一列任务保真分区，使相对状态数与最优最坏行闭合误差同时趋零；这也等价于某列趋零容限下的最小状态数为 $o(|B_N|)$。
  \item \textbf{判据的成立与失败。} 可交换故障--修复系统把 $2^N$ 个配置精确压缩为且最少需要 $N+1$ 个故障数状态，故 $\Gamma_N\to0$；带单点任务的循环链却有 $\Gamma_N=1$。这两个系统族分别说明判据可以由具体结构非平凡地满足，也绝非有限性的自动后果。
\end{enumerate}

这里的“最简”始终只表示\emph{状态数最少}。它不自动表示公式最短、参数最少、求值最快、样本最省或语义最清晰。

\subsection{与既有工作的关系}

给定分区后，有限 Markov 链的精确聚合属于 lumpability 的经典问题\citep{KemenySnell1976,Buchholz1994}；Markov functions 与线性不变性给出相关的投影条件\citep{RogersPitman1981,GurvitsLedoux2005}。概率互模拟和状态空间最小化从行为等价与分区细化角度研究相近结构\citep{LarsenSkou1991,DerisaviEtAl2003}。计算力学中的因果状态则以未来预测等价构造状态\citep{CrutchfieldYoung1989,ShaliziCrutchfield2001}。本文不把这些已有结果重新命名为“涌现定理”；本文的组织重点是把操作性复杂度基线、有限系统的精确端点与近似前沿、以及系统族的压缩--闭合判据连成一条可审计的逻辑链。

残余对称和 Transformer 预测状态可以构造候选表示，但它们不是主判据成立的必要条件。Transformer 的完整分布加权定理及其与最坏情形判据的对齐条件见附录~\ref{app:transformer-predictive-states}。Markov category 只为随机映射及其复合提供组织语言\citep{Fritz2020}；标准 Borel 空间、随机粗粒化和连续时间也各自需要附加条件。这些内容均完整保留在附录，而不被写成有限主定理已经证明的无条件推广。

\subsection{论证路线与结论边界}

正文按以下顺序推进：

\[
\begin{aligned}
&\text{操作性规范化}
\longrightarrow\text{最粗精确任务商}
\longrightarrow\text{复杂度--闭合前沿}
\\[-0.2em]
&\hspace{5em}\longrightarrow\text{纤维预测判据}
\longrightarrow\Gamma_N\text{ 系统族判据}
\longrightarrow\text{正反系统族}.
\end{aligned}
\]

第~\ref{sec:operational}节固定有意义的微观复杂度分母；第~\ref{sec:closure}节处理单个有限系统的精确端点和近似前沿；第~\ref{sec:criterion}节给出系统族判据；第~\ref{sec:fault-repair}节用一正一反两个系统族检验其非平凡性。判据完成以后，第~\ref{sec:construction}节才讨论候选表示如何由任务充分性、残余对称或预测状态产生；第~\ref{sec:verification}节说明有限检验、估计稳健性与时间边界；第~\ref{sec:discussion}节总结扩展接口和开放问题。所有独立证明、辅助反例及扩展材料统一置于附录。

本文的闭合专指给定动力学族与误差范数下的一步 Markov 闭合。结论不自动蕴含长时间准确、跨分布泛化、因果干预自治、内部神经表征、低维流形或高效发现；任何此类加强都需要另加定义和条件。
